\begin{lemma}[Exceptional-row restoration]
Let $H,F$ be finite simple graphs on the same vertex set, with
disjoint covering parts $U,W$ of size $n$, and with every edge of
$H$ crossing those parts. Let $U_0\subseteq U,W_0\subseteq W$.
Let real numbers $c,h,\eta$ satisfy
\[
0<c<1,\quad0<h\leq c/100,\quad0<\eta<1,\quad
3h<\eta,\quad4h<\eta,\quad2h+4h/\eta<\eta,\quad hn\geq1.
\]
Assume $|U\setminus U_0|,|W\setminus W_0|\leq2hn$;
every vertex of $H$ has at least $7cn/4$ neighbors in the opposite
full part; $F\subseteq H$ has edges only between $U_0,W_0$;
the opposite-core degrees of $F$ lie in $[(c-4h)n,(c+2h)n]$;
and $|d_F(A,B)-c|\leq2h$ whenever
$A\subseteq U_0,B\subseteq W_0$ and $|A|,|B|>hn$.
There is a graph $G'\subseteq H$ on the same vertex set such that
$(U,W)$ is $(\eta,c)$-super-regular, with the parameters regarded
as extended nonnegative reals.
\end{lemma}
\begin{proof}
Write $B_U=U\setminus U_0$, $B_W=W\setminus W_0$.
Since $h>0$ and $hn\geq1$, $n>0$ and $1/n\leq h$.
Let $k=\lfloor cn\rfloor$, interpreted as a natural number since
$cn>0$. Then $cn-1\leq k\leq cn$.
For any exceptional vertex in $B_U$, deletion of $B_W$ from
its possible neighbors loses at most $2hn$, so it has at least
$(7c/4-2h)n\geq cn$ neighbors in $W_0$. The last inequality
uses $2h\leq c/50<3c/4$. Likewise every vertex of $B_W$
has at least $cn$ neighbors in $U_0$. Each of these finite
neighbor sets therefore has cardinality at least the integer $k$.
Choose a $k$-element subset for each exceptional vertex.
Define $G'$ by adding to $F$ exactly the edges from these
exceptional vertices to their chosen opposite-core neighbors,
in both orientations. All are edges of $H$. No added edge has
two exceptional endpoints, and no edge is prescribed by two
exceptional vertices. This is a simple graph on the original
vertex set, contained in $H$, with decidable adjacency.

Every exceptional vertex has degree exactly $k$. Every core
vertex gains at most $|B_W|\leq2hn$ neighbors if it is in $U_0$,
or at most $|B_U|\leq2hn$ if it is in $W_0$.
The support condition on $F$ makes its degree into the full
opposite part equal its degree into the opposite core.
Thus core degrees in $G'$ lie in $[(c-4h)n,(c+4h)n]$.
For exceptional degrees use
$k\geq cn-1\geq(c-h)n\geq(c-4h)n$ and
$k\leq cn\leq(c+4h)n$.
Since $4h<\eta$, all opposite-part degrees therefore lie in
$[(c-\eta)n,(c+\eta)n]$ as real inequalities.

For a regularity test take $X\subseteq U,Y\subseteq W$ with
$a=|X|>\eta n$, $b=|Y|>\eta n$. These imply $a,b>0$.
Put $X_0=X\cap U_0,Y_0=Y\cap W_0$. The complement size
bounds and $3h<\eta$ give
$|X_0|\geq a-2hn>hn$, and likewise $|Y_0|>hn$.
The number $R=ab-|X_0||Y_0|$ of tested positions with at least
one exceptional endpoint satisfies
\[
0\leq R\leq|X\cap B_U|b+a|Y\cap B_W|
\leq2hn(a+b)<(4h/\eta)ab.
\]
For the strict last inequality, use $n/a<1/\eta$ and
$n/b<1/\eta$, valid since $n,\eta,a,b>0$.
On the other positions $G'$ agrees with $F$.
The assumed core density bound, after multiplying by
$|X_0||Y_0|>0$, gives
\[
(c-2h)|X_0||Y_0|\leq e_F(X_0,Y_0)
\leq(c+2h)|X_0||Y_0|.
\]
Here $0<c-2h<1$, by $h\leq c/100$ and $c<1$.
At most $R$ added edges occur in $X\times Y$. Consequently
\[
\begin{split}
e_{G'}(X,Y)&\geq(c-2h)(ab-R)
             \geq(c-2h-4h/\eta)ab,\\
e_{G'}(X,Y)&\leq(c+2h)|X_0||Y_0|+R
             \leq(c+2h+4h/\eta)ab.
\end{split}
\]
The budget $2h+4h/\eta<\eta$ and division by $ab>0$ prove
the lower and upper real density inequalities with error $\eta$.

It remains to verify precisely the extended-real predicates.
By \noderef{PairDensityReal}, on these nonempty tested sets the
extended density is finite and its real value equals the density
just bounded. All sizes, degrees, and upper endpoints are finite,
so their real upper inequalities translate directly to extended
upper inequalities. For every lower endpoint, if $c\geq\eta$,
the extended truncated difference has real value $c-\eta$ and
the real lower inequality translates directly, including its
multiplication by the opposite-part size in the degree clauses.
If $c<\eta$, the truncated difference is zero and the extended
lower bounds follow from nonnegativity, without treating a
negative real number as an extended parameter. The size tests
translate exactly as well, since $\eta$ and all sizes are finite.
Thus the density tests are those of \noderef{RegularPair}, and
the two degree clauses are those of \noderef{SuperRegularPair}.
The graph is on the unchanged vertex set throughout, proving the
claimed contained spanning conclusion.
\end{proof}
